% !TEX TS-program = xelatex
% !TEX encoding = UTF-8 Unicode
% !Mode:: "TeX:UTF-8"

\documentclass{resume_en}
% \usepackage{zh_CN-Adobefonts_external} % Simplified Chinese Support using external fonts (./fonts/zh_CN-Adobe/)
% \usepackage{zh_CN-Adobefonts_internal} % Simplified Chinese Support using system fonts
\usepackage{linespacing_fix} % disable extra space before next section
\usepackage{cite}

\begin{document}
\pagenumbering{gobble} % suppress displaying page number

\name{Minghan Wang}

% {E-mail}{mobilephone}{homepage}
% be careful of _ in emaill address
\contactInfo{(+61) 421-098-609}{wangmh1994@gmail.com}{PhD@Monash University}{GitHub@yuriak}
% {E-mail}{mobilephone}
% keep the last empty braces!
%\contactInfo{xxx@yuanbin.me}{(+86) 131-221-87xxx}{}
 
% \section{个人总结}
% 本人乐于钻研，工作勤奋，热爱尝试挑战性的工作，痴迷于机器学习算法的研究
% 本人在校成绩优秀、乐观向上，工作负责、自我驱动力强、热爱尝试新事物，认同开放、连接、共享的Web在未来的不可替代性。在校期间长期从事可视分析(Web的2D/3D时空可视化)相关研究，对Web技术发展趋势及前端工程化解决方案有浓厚兴趣。\textbf{现任职于华为2012实验室。}

% \section{\faGraduationCap\ 教育背景}
\section{Working Experience}
\datedsubsection{\textbf{Huawei 2012 Lab --- Text Machine Translation Lab}\ \ \ \  \textit{Algorithm Engineer}}{2020.1-2023.6}
Conducting research, publishing papers, and participating in algorithm competitions to improve the influence of the team. Meanwhile, taking part in the development of Huawei's self-developed machine translation engine (Huawei Translate) specifically on features of speech translation and OCR translation, which are embedded in the Harmony OS and HMS Core.
\begin{itemize}
  \item Currently working on developing the streaming ASR and ST system used for subtitles translation.
  \item Developed and optimized the OCR algorithm in the OCR translation feature of the Huawei Translate product, obtained SOTA performance (ACC>93.7\%).
  \item Developed APE engine to support translators in the department.
  \item Participated in 2020/2021/2022 WMT and IWSLT competitions, obtained first and second places.
  \item Published more than 5 papers to improve the research ability and influence of the team.
  \item Led the Non-autoregressive Translation and Speech Translation research projects in the lab.
\end{itemize}

\datedsubsection{\textbf{Huawei 2012 Lab --- RAMS Lab}\ \ \ \ \textit{Algorithm Engineer} }{2019.1-2020.1}
Constructing the lightweight AI customer service system used within the company, opening up inter-departmental data retrieval channels, and improving the efficiency of inter-departmental information access.
\begin{itemize}
  \item Developed the intent recognition and slot filling algorithms for low-resource scenarios, improved the accuracy under limited training data.
  \item Unified the intent recognition and slot filling components of two QA systems (IR-based and KG-based).
  \item Optimized the performance and efficiency of NLP components in existing systems.
\end{itemize}


\section{Education Experience}
\datedsubsection{\textbf{University of Monash}\\ Doctor of Philosophy in Computer Science}{2023.7 - Present}
\datedsubsection{\textbf{University of Melbourne}\\ Master of Science in Computer Science}{2017.1 - 2019.1}
\begin{itemize}
    \item GPA:3.75/4.0, \textbf{Top 15\% ranking, with distinction}
    \item Courses: Machine Learning, NLP, AI and Planing, Distributed System, Cloud Computing etc.
    \item Thesis: Learning Algorithms in Quantitative Trading (supervised by \textbf{Trevor Cohn})
    % \item Research Discipline: Applying Reinforcement Learning and NLP in Quantitative Trading. A task-specified RNN agent is proposed to learn to trade a portfolio with policy gradient. Further improvement (70+\% of annual return) can be obtained in the experiment of VIX futures backtest trading when introducing news titles as the feature.
\end{itemize}

\datedsubsection{\textbf{Dongbei University of Finance and Economics} \\Bachelor of Management in Information Management and Information Systems}{2012.9 - 2016.7}
\begin{itemize}
    \item GPA:3.0/4.0, \textbf{Top 20\% ranking}
    \item Thesis: The development of Online Evaluation System (supervised by \textbf{Hong Li})
    \item Served as a research assistant for the National Education Fund Project and led the development of two teaching assistant systems to improve teaching efficiency and collect lecture data.
    \item Served as a research assistant in the big data research group, participated in the research of text analysis in social networks, developed systems with SVM, Naive Bayes to conduct public opinion monitoring.
\end{itemize}

\section{Project}
% \datedsubsection{\textbf{RLQuant}}{2018.4 - 2019.1}
% \begin{itemize}
%     \item 将强化学习应用于量化交易的研究 \footnote{https://github.com/yuriak/RLQuant} (研究生毕业论文项目)
%     \item 模型通过在回测中探索从而学会做出正确的交易决策，最大化长期收益
%     \item 在优化agent的网络结构并加入一个有监督任务稳定训练后，性能可超越原始策略梯度算法10\%
%     \item 在模型中进一步引入新闻标题后可在VIX期货交易回测中获得70+\%的年化收益
% \end{itemize}

\datedsubsection{\textbf{DLQuant}}{2018.4 - 2019.1}
\begin{itemize}
    \item Applying deep learning and NLP in quantitative trading. \footnote{https://github.com/yuriak/DLQuant} (research project of master thesis)
    \item Performing price prediction of stocks and futures based on news titles.
    \item The model is able to learn the relevance of news and price movements with an attention based LSTM network.
    \item The performance is competitive to SOTA models at the time, but is more interpretable.
\end{itemize}

\datedsubsection{\textbf{RLQuant}}{2018.4 - 2019.1}
\begin{itemize}
    \item Applying reinforcement learning in quantitative trading \footnote{https://github.com/yuriak/RLQuant} (research project of master thesis)
    \item Modeling the trading process with reinforcement learning to maximize the long term return in the market.
    \item A policy gradient based model is proposed but with better stability and performance.
    \item When integrating news features, the model can obtain 70\%+ annual return in the VIX futures trading backtest.
\end{itemize}

\datedsubsection{\textbf{Bicycle racing venue management system}}{2016.8 - 2016.11}
\begin{itemize}
    \item Internship as a full-stack engineer in Shenzhen Wedo Sports.
    \item The system contains a committee scheduling panel and an APP installed on staffs' PDAs for RFID scanning.
    \item It has been deployed and served for several racing event with more than 1500 participants.
\end{itemize}

\section{Research Skills}
% increase linespacing [parsep=0.5ex]
\begin{itemize}[parsep=0.2ex]
  \item \textbf{Language \& Frameworks}: Python, Shell, Java, C/C++, C\#, JavaScript, HTML/CSS, Haskell
  \item \textbf{Research Interests}: Non-autoregressive Translation/Generation, Offline and Simultaneous Speech Recognition/Translation, Refinforcement Learning, Multi-modal Tasks, Dialogue System, Quantitative Trading
\end{itemize}

\section{Awards}
\begin{itemize}[parsep=0.2ex]
  \item First place of Huawei's Mindspore algorithm competition
  \item The 2020 annual quality award in Huawei
  \item First and second places in IWSLT 2022 Offline/Simultaneous/Speech to Speech/Isometirc Speech Translation
  \item First place in IWSLT 2021 Offline Speech Translation
  \item First place in IWSLT 2020 Video Speech Translation
  \item First places in WMT 2021 Domain, APE, QE, News and several other tracks
  \item First and second places in WMT 2020 QE, APE, Domain and news
  \item Second place in SemEval-2019 toponym resolution competition
\end{itemize}

\section{Publications}
\begin{itemize}[parsep=0.2ex]
  \item Zephyr: Zero-Shot Punctuation Restoration (ICASSP 2023)
  \item WhiSLU: End-to-End Spoken Language Understanding with Whisper (Interspech 2023)
  \item Diformer: Directional Transformer for Neural Machine Translation (EAMT 2022)
  \item Capture Human Disagreement Distributions by Calibrated Networks for Natural Language Inference (ACL 2022 finding)
%   \item Joint-training on Symbiosis Networks for Deep Nueral Machine Translation models (ICON 2021)
%   \item Self-Distillation Mixup Training for Non-autoregressive Neural Machine Translation (ICON 2021)
  \item How does Length Prediction Influence the Performance of Non-Autoregressive Translation? (BlackboxNLP 2021)
  \item HI-CMLM: Improve CMLM with Hybrid Decoder Input (INLG 2021)
  \item Make the Blind Translator See The World: A Novel Transfer Learning Solution for Multimodal Machine Translation (AMTA 2021)
  \item Efficient Transfer Learning for Quality Estimation with Bottleneck Adapter Layer (EAMT 2020)
  \item Unified Humor Detection Based on Sentence-pair Augmentation and Transfer Learning (EAMT 2020)
  \item SMAC: An Interpretable Reasoning Network for Visual Question Answering (RCQA 2020)
  \item The HW-TSC's Offline Speech Translation Systems for IWSLT 2021 Evaluation (IWSLT 2021)
  \item The HW-TSC Video Speech Translation System at IWSLT 2020 (IWSLT 2020)
  \item HW-TSC's Participation at WMT 2020 Quality Estimation Shared Task (WMT 2020)
  \item HW-TSC's Participation at WMT 2020 Automatic Post Editing Shared Task (WMT 2020)
  \item Huawei's Submissions to the WMT20 Biomedical Translation Task (WMT 2020)
  \item HW-TSC's Participation in the WMT 2020 News Translation Shared Task (WMT 2020)
  \item UniMelb at SemEval-2019 Task 12: Multi-model Combination for Toponym Resolution (SemEval 2019)
\end{itemize}


% \ GPA:3.7/4, \textbf{排名前15\%}
% 毕业论文: Online Evaluation System的设计与开发

% \section{\faCogs\ IT 技能}


% \end{itemize}

% \section{实习经历}
% \datedsubsection{\textbf{阿里巴巴集团 | Alibaba}, 前端开发工程师}{2017.6-2017.9}
% \begin{itemize}
% %   \item 飞猪北京前端团队全面负责各交通线的票务(机票/火车票/汽车票) web 应用与事业群基础架构研发
%   \item 独立负责车站地图开发(React),通过HTML5 本地存储及JSBridge实现在阿里全系应用中发布上线
%   \item 独立负责BU SPM chrome插件开发,支付成功/订单详情等页面的开发与交叉营销的接入工作
% \end{itemize}

% \datedsubsection{\textbf{北京腾云天下科技有限公司 | TalkingData},数据挖掘与可视化工程师}{2015.11-2017.5}
% \begin{itemize}
%   \item \textbf{利用海量用户定位数据，对城市空间及人群移动特征进行研究。}第一个课题是基于香农熵和人群出行模式，构建城市网格与用户矩阵分析城市多样性/流动性分布；可视分析平台前端与可视化基于D3/Vue/Express开发，数据分析与存储采用Python/MySQL/MongoDB技术，为了均衡大数据情况下的页面可视化渲染消耗用canvas替代svg。第二个课题是对海量商场定位数据做人群分类与可视化查询，依据该系统撰写的论文被CIKM 2016(DAVA Workshop)录用，并收录于中科院软件所年会成果集
%   \item 负责数据科学部HQ LAB的可视化原型开发，主导 TalkingMind 平台系统设计与前端开发
% \end{itemize}

% \datedsubsection{\textbf{北京格灵深瞳信息技术有限公司 | DeepGlint},Web开发工程师}{2015.7-2015.9}
% \begin{itemize}
%   \item \textbf{独立负责MUSE部门的可视化组件研发。}与平台研发、设计协作完成 DeepGlint Developer 平台可视化图表组件的集成开发，符合完全定制化渲染、响应式布局与实时更新等特点
%   \item 利用 D3+Vue+WebGL(Three.js) 尝试实现三维空间的人群移动可视化
% \end{itemize}

% % \begin{onehalfspacing}
% % \end{onehalfspacing}

% % \datedsubsection{\textbf{DID-ACTE} 荷兰莱顿}{2015年3月 - 2015年6月}
% % \role{本科毕业设计}{LIACS 交换生}
% % 利用结巴分词对中国古文进行分词与词性标注，用已有领域知识训练形成 classifier 并对结果进行调优
% % \begin{onehalfspacing}
% % \begin{itemize}
% %   \item 利用结巴分词对中国古文进行分词与词性标注
% %   \item 利用已有领域知识训练形成 classifier, 并用分词结果进行测试反馈
% %   \item 尝试不同规则，对 classifier 进行调优
% % \end{itemize}
% % \end{onehalfspacing}

% \section{竞赛获奖/项目作品}
% % increase linespacing [parsep=0.5ex]
% \begin{itemize}[parsep=0.2ex]
% %   \item LeetCodeOJ Solutions, \textit{https://github.com/hijiangtao/LeetCodeOJ}
%   \item 第三届中国软件杯大学生软件设计大赛\textbf{全国一等奖}( \textit{http://www.cnsoftbei.com/} ),2014 年8月
%   \item 中国机器人大赛创意设计大赛\textbf{全国特等奖}( \textit{http://www.rcccaa.org/} ),2013年8月
% %   \item 中国机器人大赛暨Robocup公开赛(武术擂台赛)全国一等奖,2013年10月
%   \item 第11届北京理工大学“世纪杯”竞赛学生课外科技作品竞赛\textbf{特等奖},2013年8月
%   \item VIS Components for security system, \textit{https://hijiangtao.github.io/ss-vis-component/}
%   \item 个人博客：\textit{https://hijiangtao.github.io/}，更多作品见 \textit{https://github.com/hijiangtao}
% %   \item 电视节目"爸爸去哪儿"可视化分析展示, \textit{https://hijiangtao.github.io/variety-show-hot-spot-vis/}
% \end{itemize}

% % \section{\faHeartO\ 项目/作品摘要}
% % \section{项目/作品摘要}
% % \datedline{\textit{An Integrated Version of Security Monitor Vis System}, https://hijiangtao.github.io/ss-vis-component/ }{}
% % \datedline{\textit{Dark-Tech}, https://github.com/hijiangtao/dark-tech/ }{}
% % \datedline{\textit{融合社交网络数据挖掘的电视节目可视化分析系统}, https://hijiangtao.github.io/variety-show-hot-spot-vis/}{}
% % \datedline{\textit{LeetCodeOJ Solutions}, https://github.com/hijiangtao/LeetCodeOJ}{}
% % \datedline{\textit{Info-Vis}, https://github.com/ISCAS-VIS/infovis-ucas}{}


% % \section{\faInfo\ 社会实践/其他}
% \section{社区参与/实践其他}
% % increase linespacing [parsep=0.5ex]
% \begin{itemize}[parsep=0.2ex]
%   \item 乐于参与开源社区讨论,\textbf{参与翻译 Vue.js, webpack, WebAssembly, Babel 文档，印记中文成员}
%   \item 中国科学院大学2016秋季学期可视化与可视分析课程助教，\textit{http://vis.ios.ac.cn/infovis-ucas/}
%   \item 未来论坛学生会成员、北理社联新闻信息中心主任、北理工软件学院学生会宣传部副部长(2012-2016)
%   \item 2013-2015 北京市共青团“温暖衣冬”志愿者,第九届园博会志愿者,2014 FLL机器人世锦赛志愿者
% \end{itemize}

%% Reference
%\newpage
%\bibliographystyle{IEEETran}
%\bibliography{mycite}
\end{document}
